
\SGARefTarget{sga3:xvi:section:4}%
\section*{4. Groups with Affine Fibers, of Unipotent Rank Zero}
\addcontentsline{toc}{section}{4. Groups with Affine Fibers,
of Unipotent Rank Zero}

\noindent\SGARefTarget{sga3:xvi:theorem:4:1}%
\textbf{Theorem 4.1.} --- \textit{Let \(S\) be a prescheme and
let \(G\) be an \(S\)-group prescheme smooth over \(S\), of finite
presentation, with affine, connected fibers, of unipotent rank
zero (\SGARefLink{sga3:xv:definition:6:1:ter}{Exposé XV, 6.1 ter}). Then:}

\begin{enumerate}
\renewcommand{\labelenumi}{\textit{\alph{enumi})}}
\item \textit{The center \(Z\) of \(G\) is a group of
multiplicative type (and is therefore a reductive center of \(G\);
\SGARefLink{sga3:xii:4.1}{Exposé XII, 4.1}).}

\item \textit{\(G\) satisfies conditions a) through d) of \SGARefLink{sga3:xvi:theorem:3:3}{3.3}.}

\item \textit{Locally for the étale topology, \(G\) possesses
maximal tori, and these are also Cartan subgroups of \(G\). The
functor \(\mathscr{TM}_G\) of maximal tori of \(G\) is
representable by an \(S\)-prescheme smooth and affine over \(S\).}

\item \textit{\(G\) is quasi-affine over \(S\). Moreover, if
\(T\) is a maximal torus of \(G\), if}
\[
  N=\Norm_G(T)
\]
\textit{is its normalizer in \(G\), and if \(W=N/T\) is the Weyl
group relative to \(T\), the following conditions are equivalent:}
\begin{enumerate}
\renewcommand{\labelenumii}{\textit{(\roman{enumii})}}
\item \textit{\(G\) is affine over \(S\);}
\item \textit{\(G'=G/Z\) is affine over \(S\);}
\item \textit{\(N\) is affine over \(S\);}
\item \textit{\(W\) is affine over \(S\).}
\end{enumerate}
\textit{These conditions are always satisfied if \(S\) is locally
noetherian of dimension at most \(1\).}
\end{enumerate}

\medskip
\noindent\textit{Proof.}
Since \(G\) has affine fibers (and hence abelian rank zero) and
unipotent rank zero, \(G\) possesses maximal tori locally for the
étale topology (\SGARefLink{sga3:xv:corollary:8-18}{Exposé XV, 8.18}), and every maximal torus of \(G\)
is plainly a Cartan subgroup of \(G\). This proves the first part
of c). To see that \(Z\) is of multiplicative type, we may suppose,
by what precedes, that \(G\) possesses a maximal torus \(T\).
Since \(T\) is also a Cartan subgroup, it contains \(Z\) (because
\(T=\Centr_G(T)\) by \SGARefLink{sga3:xii:7.1}{Exposé XII, 7.1}), and
\[
  Z=Z\cap T
\]
is of multiplicative type by \SGARefLink{sga3:xvi:proposition:3:5}{3.5} i). This proves a). Assertion b)
is clear in view of a). On the other hand, the functor
\(\mathscr{TM}_G\) is isomorphic to the functor of Cartan
subgroups of \(G\) (\SGARefLink{sga3:xii:7.1}{Exposé XII, 7.1}), and is therefore smooth and
affine over \(S\) by \SGARefLink{sga3:xvi:theorem:3:3}{3.3} c). This finishes the proof of c). It
remains to prove d).

\medskip
\noindent\textit{Proof of d).}
Since \(G'\) is quasi-affine (\SGARefLink{sga3:xvi:theorem:3:3}{3.3} b)) and \(Z\) is of
multiplicative type, hence affine over \(S\), \(G\) is quasi-affine
(Exposé VI\textsubscript{B}, \S9).

\medskip
\noindent i) \(\Longleftrightarrow\) ii).
If \(G\) is affine, then \(G'\) is affine by \SGARefLink{sga3:ix:proposition:2:3}{Exposé IX, 2.3}. If
\(G'\) is affine, then \(G\) is affine as an extension of an affine
group by an affine group (Exposé VI\textsubscript{B}, \S9).

\medskip
\noindent i) \(\Longleftrightarrow\) iii).
If \(G\) is affine, then \(N\) is affine, since it is closed in
\(G\) (\SGARefLink{sga3:xi:corollary:6:11}{Exposé XI, 6.11} a)). Moreover, \(G/N\) is isomorphic to the
functor \(\mathscr{TM}_G\) (conjugacy of maximal tori; cf.\
\SGARefLink{sga3:xii:7.1}{Exposé XII, 7.1} b)), and is therefore affine over \(S\) by c).
Thus if \(N\) is affine over \(S\), then \(G\) is affine, as an
``extension'' of an affine homogeneous space by an affine group
(Exposé VI\textsubscript{B}, \S9).

\medskip
\noindent iii) \(\Longleftrightarrow\) iv).
If \(N\) is affine, then \(W=N/T\) is affine (\SGARefLink{sga3:ix:proposition:2:3}{Exposé IX, 2.3}), and
the converse follows from Exposé VI\textsubscript{B}, \S9.

Moreover, \(N/T\) is representable by an \(S\)-group prescheme
that is étale, separated, and of finite type over \(S\) (\SGARefLink{sga3:xv:proposition:7-1}{Exposé XV,
7.1} iv)), hence quasi-finite over \(S\). The last assertion of d)
therefore follows from the next lemma, applied with \(X=W\) and
\(Y=S\).

\medskip
\noindent\SGARefTarget{sga3:xvi:lemma:4:2}%
\textbf{Lemma 4.2.} --- \textit{Let \(S\) be a locally noetherian
prescheme of dimension \(1\), let \(X\) and \(Y\) be two
\(S\)-preschemes locally of finite type over \(S\), and let
\(u:X\to Y\) be a quasi-finite and separated \(S\)-morphism.
Suppose that, for every point \(s\) of \(S\), the morphism}
\[
  u_s:X_s\longrightarrow Y_s
\]
\textit{induced from \(u\) by the base change
\(\operatorname{Spec}\kappa(s)\to S\) is finite. Then \(u\) is an
affine morphism.}

\medskip
The assertion to be proved is local on \(Y\), which allows us to
suppose that \(Y\) (and hence also \(X\)) is of finite type over
\(S\). By EGA IV, \S8, it suffices to prove 4.2 when \(S\) is the
spectrum of a local ring \(A\). By fpqc descent we may suppose
that \(A\) is complete, then that \(A\) is reduced (EGA II,
1.6.4), and then that \(A\) is normal (Nagata's theorem, EGA
0\textsubscript{IV}, \S22, and Chevalley's theorem, EGA II,
6.7.1). If \(A\) is a field, then \(u\) is finite by hypothesis,
hence affine. Otherwise \(A\) is a discrete valuation ring; let
\(s\) and \(t\) be the closed and generic points of \(S\), and let
\(\pi\) be a uniformizer of \(A\). Let \(y\) be a point of \(Y_s\).
Applying EGA IV, \S8, once more, we may replace \(Y\) by the
spectrum \(Y'\) of \(\mathcal O_{Y,y}\), and \(X\) by
\[
  X'=X\times_Y Y';
\]
we may even suppose that \(\mathcal O_{Y,y}\) is complete. Since
\(u\) is quasi-finite and separated, EGA II, 6.2.6, shows that
\(X'\) is the disjoint union of two schemes \(X_1\) and \(X_2\),
where \(X_1\) is finite over \(Y'\) (hence affine), while
\(u(X_2)\) does not contain the closed point \(y\) of \(Y'\). I
claim that, under these conditions, \(X_2\) does not meet the
closed fiber \(X'_s\). Indeed, by hypothesis,
\[
  X'_s\longrightarrow Y'_s
\]
is finite, so the restriction of this morphism to the closed
subset \((X_2)_s\) is finite, and its image in \((Y')_s\) is
closed. Since this image does not contain the closed point of the
local scheme \((Y')_s\), \((X_2)_s\) is empty. Since \(u_t\) is
finite, the restriction to \((X_2)_t=X_2\) of the morphism
\[
  X'_t\longrightarrow Y'_t
\]
is finite. Moreover \(Y'_t=Y'_\pi\), so the open immersion
\(Y'_t\to Y'\) is affine. In short, the composite morphism
\[
  X_2\longrightarrow Y'_t\longrightarrow Y'
\]
is affine, and it follows that
\[
  X'=X_1\amalg X_2\longrightarrow Y'
\]
is affine.

\medskip
\noindent\SGARefTarget{sga3:xvi:corollary:4:3}%
\textbf{Corollary 4.3.} --- \textit{Let \(S\) be a prescheme, and
let \(G\) be an \(S\)-group prescheme smooth and of finite
presentation over \(S\), with affine, solvable, connected fibers,
of unipotent rank zero. Then \(G\) is affine over \(S\). If,
moreover, the center of \(G\) is the unit group and \(S\) is
quasi-compact, the canonical morphism}
\[
  i_n:G\longrightarrow\operatorname{GL}_S(G^{(n)})
  \qquad\textit{(cf.\ 3.1 b))}
\]
\textit{is a closed immersion for all sufficiently large \(n\).}

\medskip
To prove that \(G\) is affine over \(S\), we may suppose that
\(G\) possesses a maximal torus \(T\) (\SGARefLink{sga3:xvi:theorem:4:1}{4.1} c)). Since \(G\) has
solvable fibers, the Weyl group relative to \(T\) is the unit
group (BIBLE, 6, Th.\ 1, Cor.\ 3), and condition \SGARefLink{sga3:xvi:theorem:4:1}{4.1} d) iv) is
satisfied. If the center of \(G\) is the unit group, \(i_n\) is a
monomorphism for all sufficiently large \(n\) (\SGARefLink{sga3:xvi:proposition:3:1}{3.1}), and is
therefore a closed immersion (\SGARefLink{sga3:xvi:corollary:1:5}{1.5} b)).
